\documentclass[11pt,a4paper]{article}

% ==============================================================================
% PACKAGES
% ==============================================================================
\usepackage[a4paper,
            top=2.3cm, bottom=2.3cm,
            left=2.3cm, right=2.3cm,
            headsep=12pt, footskip=12pt]{geometry}
\usepackage{graphicx}
\usepackage{booktabs}
\usepackage{amsmath, amssymb}
\usepackage{siunitx}
\usepackage{caption}
\usepackage{float}
\usepackage{fancyhdr}
\usepackage{hyperref}
\usepackage{cleveref}
\captionsetup{font=small, labelfont=bf}

% Clean header/footer with no extra rules
\pagestyle{fancy}
\fancyhf{}
\fancyfoot[C]{\thepage}
\renewcommand{\headrulewidth}{0pt}
\renewcommand{\footrulewidth}{0pt}

\graphicspath{{D:/SoSe 2026/MHYWI05 Statistical Learning for Earth System Sciences/Tasks for the report/Figures/}}

% Tighter float spacing (this was already in your working version)
\setlength{\textfloatsep}{6pt plus 2pt minus 2pt}
\setlength{\floatsep}{6pt plus 2pt minus 2pt}
\setlength{\intextsep}{6pt plus 2pt minus 2pt}

% ==============================================================================
\begin{document}
% ==============================================================================

% ------------------------------------------------------------------------------
% TITLE PAGE
% ------------------------------------------------------------------------------
\pagenumbering{gobble}
\thispagestyle{empty}
\begin{titlepage}
    \centering
    
    % --- TOP: LOGO (alone, left-aligned) ---
    \begin{flushleft}
        \includegraphics[width=0.5\textwidth]{TUDlogo.png}
    \end{flushleft}
    
    \vspace{0.1cm} % Space between logo and text block
    
    % --- TEXT BLOCK (below logo, left-aligned) ---
    \begin{flushleft}
        \footnotesize
        % First dash line
        \rule{\textwidth}{0.1pt} \\[0.1cm]
        % Text between the two dash lines
        Faculty of Environmental Sciences, Department of Hydro Sciences,\\
        Institute of Hydrology and Meteorology
    \end{flushleft}
    
    \vspace{1.5cm} % Space before module description
    
    % --- CENTERED: Module description ---
    {\Large Project Report in the Module\\
     MHYWI05 ``Statistical Learning for Earth System Sciences''\par}
    
    \vspace{1.5cm} % Space before Title
    
    % --- CENTERED: Title ---
    {\Large \textbf{Statistical Learning for Earth System Sciences:\\
     Crop Yield Prediction and European Climate Trends}\par}
    
    \vspace{2.5cm} % Space before metadata
    
    % --- BOTTOM: Left-aligned Author Info Table ---
    \raggedright
    \begin{tabular}{@{}p{4.5cm} p{9.5cm}@{}}
        \textbf{Module Code:} & MHYWI05 \\
        \vspace{0.7cm} \\
        \textbf{Study Program:} & M.Sc. Hydro Science and Engineering \\
        \vspace{0.7cm} \\
        \textbf{Student:} & Galib Ahmed Chowdhury (5035774) \\
        \vspace{0.7cm} \\
        \textbf{Responsible Lecturer:} & Prof. Dr. Jakob Zscheischler \\
        \vspace{0.3cm} \\
         & \texttt{jakob.zscheischler@tu-dresden.de} \\
        \vspace{0.7cm} \\
        \textbf{Institution:} & Faculty of Environmental Sciences, TU Dresden \\
        \vspace{0.3cm} \\
         & Institute of Hydrology and Meteorology \\
        \vspace{0.7cm} \\
        \textbf{Submission Date:} & 30. September 2026 \\
    \end{tabular}
    
    \vfill
\end{titlepage}
\clearpage

% ------------------------------------------------------------------------------
% MAIN BODY
% ------------------------------------------------------------------------------
\pagenumbering{arabic}
\setcounter{page}{1}

% ==============================================================================
\section{Introduction}
% ==============================================================================

This report applies statistical learning methods to two environmental datasets
with different structures and inferential goals. The first is a panel of Iowa
crop yields (25 locations, 60 years) with monthly growing-season meteorological
predictors; here the question is predictive, and specifically how the standard
random train/test split compares to one that respects the data's temporal
structure. The second is gridded European climate data over 1950--2025; here the
question is inferential, namely how strong the long-term temperature and
precipitation trends are and how correcting for thousands of simultaneous tests
affects which trends survive. Both analyses turn on assumptions that are easy to
ignore: independence of observations in the crop panel, and independence of
tests in the climate grids.

% ==============================================================================
\section{Iowa Crop Yield Data}
% ==============================================================================

\subsection{Data and dimensionality: Answer 1a}

The dataset contains $n = 1500$ observations (25 locations $\times$ 60 years) and
$p = 23$ predictors: six monthly values each of precipitation, radiation, and
temperature during the growing season, plus five extreme-event indicators (frost
days, warm days, maximum and minimum temperature, and maximum 5-day
precipitation). There are no missing values.

PCA on the standardised predictors shows that \textbf{11 components are required
to explain at least 80\% of the total variance} (81.3\% cumulative;
Figure~\ref{fig:pca}). The scree plot shows no sharp elbow: variance declines
gradually rather than concentrating in a few components, indicating moderate but
not extreme redundancy among the predictors.

\begin{figure}[htbp]
  \centering
  \includegraphics[width=0.95\textwidth]{fig1_pca.png}
  \caption{Variance explained per principal component (left) and cumulative
  variance (right), with the 80\% threshold and the 11th component marked in red.}
  \label{fig:pca}
\end{figure}

\subsection{Violation of the i.i.d.\ assumption: Answer 1b}

The data violate the independent and identically distributed assumption through
strong cross-sectional (spatial) dependence. A representative monthly
temperature predictor, \texttt{temp\_3}, is correlated at $r = 0.958$ across the
25 Iowa locations within the same year, and yield itself shows a mean
cross-location correlation of $r = 0.685$. Growing conditions are shared across
the state, so observations from the same year are not independent draws.

A fixed-effects regression confirms this: fitting yield on location and year as
categorical factors alone, with no weather information, explains \textbf{71.8\%
of the variance}. Temporal dependence is comparatively weak --- the mean
within-location correlation between year and yield is $r = -0.094$ --- so the
identical-distribution assumption is a smaller concern than independence.

\subsection{Random train/test split: Answer 1c}

The data were split 50/50 ($n = 750$ each) using \texttt{set.seed(99)}. Two
models were fit: best subset regression with size chosen by minimum BIC, and
lasso with 5-fold cross-validation. For the lasso, both $\lambda_{\min}$ and
$\lambda_{1\mathrm{se}}$ are reported. Results are shown in Table~\ref{tab:1c}.

\begin{table}[htbp]
  \centering
  \small
  \begin{tabular}{l c S[table-format=1.4] S[table-format=1.4] S[table-format=1.4] S[table-format=1.4]}
    \toprule
    Model & Vars & {Train MSE} & {Train $R^2$} & {Test MSE} & {Test $R^2$} \\
    \midrule
    Best subset (BIC, $k=11$) & 11 & 2.8381 & 0.6704 & 3.2581 & 0.6343 \\
    Lasso $\lambda_{\min}$    & 20 & 2.7716 & 0.6782 & 3.2366 & 0.6367 \\
    Lasso $\lambda_{1\mathrm{se}}$ & 12 & 3.0566 & 0.6451 & 3.3540 & 0.6235 \\
    \bottomrule
  \end{tabular}
  \caption{Random 50/50 split ($n_{\text{train}} = n_{\text{test}} = 750$).
  Test $R^2$ is computed against the test set's own mean.}
  \label{tab:1c}
\end{table}

$\lambda_{1\mathrm{se}}$ retains 12 variables against $\lambda_{\min}$'s 20, at
a cost of only ${\sim}0.013$ in test $R^2$; it is the more defensible choice for
a report where interpretability matters.

Best subset (11 variables) selected: \texttt{pr\_2, pr\_3, rad\_2, rad\_3,
temp\_2, temp\_3, temp\_4, temp\_6, warm\_days, max\_temp, max\_5day\_pr}.

Lasso at $\lambda_{1\mathrm{se}}$ (12 variables) selected a very similar set:
\texttt{pr\_2, pr\_3, rad\_2, rad\_3, rad\_5, temp\_2, temp\_3, temp\_4,
temp\_6, warm\_days, max\_temp, max\_5day\_pr}.

The two sets differ only by \texttt{rad\_5}, which lasso includes. Test $R^2$
for both methods lands in the 0.62--0.64 range, and the
predicted-versus-observed plots in Figure~\ref{fig:1c} track the 1:1 line
closely across the full yield range.


\begin{figure}[H]
  \centering
  \includegraphics[width=0.85\textwidth]{fig2_pred_vs_obs.png}
  \caption{Predicted vs.\ observed yield on the test set (random split),
  best subset (left) and lasso $\lambda_{\min}$ (right); dashed line is 1:1.}
  \label{fig:1c}
\end{figure}

\subsection{Year-parity split and the effect of spatial leakage: Answer 1d}

The models were refitted with an alternative split: even years for training, odd
years for testing. This yields 750 observations in each set, but with no year
appearing in both training and test data. Results are shown in
Table~\ref{tab:1d}.

\begin{table}[H]
  \centering
  \small
  \begin{tabular}{l c S[table-format=1.4] S[table-format=1.4] S[table-format=1.4] S[table-format=1.4]}
    \toprule
    Model & Vars & {Train MSE} & {Train $R^2$} & {Test MSE} & {Test $R^2$} \\
    \midrule
    Best subset (BIC, $k=15$) & 15 & 2.7494 & 0.6934 & 4.5750 & 0.4648 \\
    Lasso $\lambda_{\min}$    & 23 & 2.6839 & 0.7007 & 4.7114 & 0.4489 \\
    Lasso $\lambda_{1\mathrm{se}}$ & 17 & 2.9036 & 0.6762 & 4.3304 & 0.4934 \\
    \bottomrule
  \end{tabular}
  \caption{Year-parity split ($n_{\text{train}} = n_{\text{test}} = 750$).}
  \label{tab:1d}
\end{table}

Training performance is essentially unchanged; test performance drops
substantially. Best subset test $R^2$ falls from 0.634 to 0.465, and lasso at
$\lambda_{\min}$ falls from 0.637 to 0.449. Figure~\ref{fig:1d} shows the cause
visually: the scatter around the 1:1 line is markedly wider, with predictions
for high-yield observations routinely falling short by 2--4\,t/ha.

\begin{figure}[H]
  \centering
  \includegraphics[width=0.85\textwidth]{fig3_pred_vs_obs_parity.png}
  \caption{As Figure~\ref{fig:1c}, for the year-parity split. Note the wider
  scatter around the 1:1 line.}
  \label{fig:1d}
\end{figure}

The mechanism is the spatial dependence documented in Answer 1b. Under the
random split, a test observation from, say, location 7 in a given year has
training-set neighbours from other locations in the same year. Because
temperature is correlated at $r = 0.96$ across locations within a year, those
neighbours carry much of the information needed to predict the test point's
yield --- the model interpolates rather than extrapolates. The year-parity split
removes this channel: no year appears in both sets, so the model must predict
under meteorological conditions it has never observed at any location. The
resulting test $R^2$ of 0.45--0.49 is the honest estimate of out-of-sample
performance.

% ==============================================================================
\section{European Climate Trends}
% ==============================================================================

\subsection{Per-cell trend analysis: Answer 2a}

The gridded climate data cover the European domain from approximately
$35^\circ$N to $75^\circ$N and $10^\circ$W to $30^\circ$E, at roughly
$0.25^\circ$ resolution, over 76 annual timesteps from 1950 to 2025. Annual mean
temperature ranges from $-6.98^\circ$C to $24.73^\circ$C; annual precipitation
sums range from 0 to 4938.4\,mm. Ocean cells are masked (NA).

For each land grid cell, a linear regression of the variable on year was fitted
and the null hypothesis $H_0$: slope $= 0$ was tested against the two-sided
alternative at $\alpha = 0.05$. This gives 11{,}791 temperature tests and
11{,}918 precipitation tests.

The strongest warming trend is $0.058^\circ$C\,yr$^{-1}$, and the strongest
cooling trend is $-0.012^\circ$C\,yr$^{-1}$. For precipitation, the strongest
increase is 17.1\,mm\,yr$^{-1}$ and the strongest decrease is
$-7.8$\,mm\,yr$^{-1}$; both are isolated single cells and likely artefacts of
the underlying interpolation over complex terrain.

Of the 11{,}791 temperature cells, \textbf{11{,}512 (97.6\%)} show a
significant trend at $\alpha = 0.05$. For precipitation, \textbf{4{,}799 of
11{,}918 (40.3\%)} are significant. Under a global null of no trends anywhere,
we would expect approximately 589.6 and 595.9 false positives respectively.
Both observed counts are far above this baseline, indicating genuine widespread
signal.

\begin{figure}[htbp]
  \centering
  \includegraphics[width=0.95\textwidth]{fig4_trend_maps.png}
  \caption{Linear trends in annual mean temperature (left, $^\circ$C/decade) and
  annual precipitation sums (right, mm/decade) over Europe, 1950--2025. Colour
  scales are symmetric around zero.}
  \label{fig:trends}
\end{figure}

Figure~\ref{fig:trends} shows the spatial distribution of the raw trends.
Warming is coherent across nearly the entire domain, with the strongest signal
over Scandinavia, the Baltic, and Central/Eastern Europe; cooling patches appear
only in isolated parts of the British Isles and the southeastern Balkans.
Precipitation trends are much patchier: increases dominate in Scotland, Norway,
and the Alps, while decreases cluster in Iberia, Italy, and the western Balkans.

\subsection{Multiple testing correction: Answer 2b}

The Benjamini--Hochberg procedure was applied separately to each variable's
$p$-values, with a self-chosen false discovery rate of $q = 0.05$. This
guarantees that, on average, no more than 5\% of the rejected null hypotheses
are false positives.

For temperature, the number of significant cells is essentially unchanged
(11{,}512 $\rightarrow$ 11{,}510; 0.02\% reduction). For precipitation, the count
falls from 4{,}799 to 3{,}838, a 20.0\% reduction. The expected number of false
discoveries among the surviving rejections is $q \times$ (number of rejections):
575.5 for temperature and 191.9 for precipitation.

\begin{figure}[htbp]
  \centering
  \includegraphics[width=0.95\textwidth]{fig5_fdr_maps.png}
  \caption{Grid cells with a statistically significant trend after
  Benjamini--Hochberg correction ($q = 0.05$).}
  \label{fig:fdr}
\end{figure}

After correction, the spatial patterns remain largely intact: temperature
warming is still near-ubiquitous, and precipitation retains its regional
structure (wetter across Scandinavia, drier across Iberia and the western
Mediterranean). The divergence in correction behaviour reflects a real
difference in signal-to-noise ratio. Temperature trends are large relative to
year-to-year variability and spatially coherent, so almost every raw-significant
cell corresponds to a genuine trend. Precipitation trends are smaller relative
to interannual variability and spatially patchier, so a meaningful fraction of
the raw-significant cells were likely noise --- which the FDR correction removes.

% ==============================================================================
\section{Conclusion}
% ==============================================================================

Both analyses make the same methodological point from different directions: the
assumptions underlying standard statistical learning procedures are frequently
violated by Earth system data, and the consequences show up in the results. In
the Iowa crop panel, ignoring spatial dependence via a random split yields a
test $R^2$ of 0.63--0.64; respecting it via a year-parity split yields
0.45--0.49. The 0.13--0.19 gap in $R^2$ is the cost of the random split's
implicit information leakage, and it would not survive in operational
forecasting, where future years are genuinely unseen. In the European climate
grids, failing to correct for multiple testing inflates the count of significant
trends by roughly 20\% for precipitation, while barely affecting temperature,
since the two variables differ in signal-to-noise ratio. In both cases, checking
whether the data's structure is compatible with the method's assumptions --- and
choosing the variant that accounts for that structure explicitly --- is what
makes the reported performance and significance estimates trustworthy.

% ==============================================================================
\end{document}
% ==============================================================================